\section*{Hoofdstuk \ref{ch:g11:igneous-rocks} --- Stollingsgesteenten}

\begin{solution}{exo:g11:igneous-rocks:1}
Een \omterm{def:g11:igneous-rocks:volcanic-plutonic}{vulkanisch gesteente} is ontstaan uit een magma dat aan het oppervlak is gestold, als
lava; een \omterm{def:g11:igneous-rocks:volcanic-plutonic}{plutonisch gesteente} is op diepte gekristalliseerd, binnenin andere gesteenten. Paren: basalt
en \omterm{def:g11:ocean-ridges:gabbro}{gabbro} (\omterm{def:g11:igneous-rocks:felsic-mafic}{mafisch}), \omterm{def:g11:igneous-rocks:rock-names}{andesiet} en \omterm{def:g11:igneous-rocks:rock-names}{dioriet} (\omterm{def:g11:igneous-rocks:felsic-mafic}{intermediair}), \omterm{def:g11:igneous-rocks:rock-names}{ryoliet} en graniet
(\omterm{def:g11:igneous-rocks:felsic-mafic}{felsisch}).
\end{solution}

\begin{solution}{exo:g11:igneous-rocks:2}
(1) Glasachtig. (2) Grofkorrelig. (3) \omterm{def:g11:igneous-rocks:porphyritic}{Porfirisch} (\omterm{def:g11:igneous-rocks:porphyritic}{fenokristen} van olivijn in een
fijne grondmassa). (4) Fijnkorrelig.
\end{solution}

\begin{solution}{exo:g11:igneous-rocks:3}
\qty{48}{\percent}: \omterm{def:g11:igneous-rocks:felsic-mafic}{mafisch}, \omterm{def:g11:ocean-ridges:gabbro}{gabbro}. \qty{60}{\percent}: \omterm{def:g11:igneous-rocks:felsic-mafic}{intermediair}, \omterm{def:g11:igneous-rocks:rock-names}{dioriet}.
\qty{72}{\percent}: \omterm{def:g11:igneous-rocks:felsic-mafic}{felsisch}, graniet.
\end{solution}

\begin{solution}{exo:g11:igneous-rocks:4}
Een \omterm{def:g11:igneous-rocks:intrusion}{gang} is een plaat stollingsgesteente die de lagen doorsnijdt van de gesteenten waarin
hij is binnengedrongen; een \omterm{def:g11:igneous-rocks:intrusion}{sill} is een plaat evenwijdig aan de lagen, ertussen; een \omterm{def:g11:igneous-rocks:intrusion}{pluton} is
een grote \omterm{def:g11:igneous-rocks:intrusion}{intrusie} van willekeurige vorm die op diepte is gevormd. De verticale \omterm{def:g11:igneous-rocks:intrusion}{gang} die
van de \omterm{def:g11:igneous-rocks:intrusion}{pluton} omhoogloopt, voedde de vulkaan.
\end{solution}

\begin{solution}{exo:g11:igneous-rocks:5}
Olivijn ontstaat als eerste; kaliveldspaat, kwarts en soms muscoviet
als laatste. Plagioklaas begint rijk aan calcium en wordt steeds rijker aan
natrium terwijl het magma afkoelt.
\end{solution}

\begin{solution}{exo:g11:igneous-rocks:6}
(1) A: \qty{74}{\percent} \ce{SiO2}, \omterm{def:g11:igneous-rocks:felsic-mafic}{felsisch}; B: \qty{47}{\percent}, \omterm{def:g11:igneous-rocks:felsic-mafic}{mafisch}.
(2) A is arm aan \ce{FeO}, \ce{MgO} en \ce{CaO} en rijk aan \ce{Na2O} en
\ce{K2O}, zoals een \omterm{def:g11:igneous-rocks:felsic-mafic}{felsisch} magma; B is rijk aan \ce{FeO}, \ce{MgO} en \ce{CaO},
arm aan \ce{Na2O}, zoals een \omterm{def:g11:igneous-rocks:felsic-mafic}{mafisch} magma. (3) A is een \omterm{def:g11:igneous-rocks:rock-names}{ryoliet}; B een \omterm{def:g11:ocean-ridges:gabbro}{gabbro}.
\end{solution}

\begin{solution}{exo:g11:igneous-rocks:7}
De tijd groeit als het kwadraat van de dikte. \qty{1}{m}: $10/4 = 2.5$
dagen. \qty{20}{m}: $10 \times 10^2 = 1000$ dagen, ongeveer 2.7 jaar.
\qty{200}{m}: $10 \times 100^2 = \num{100000}$ dagen, ongeveer 270 jaar.
\end{solution}

\begin{solution}{exo:g11:igneous-rocks:8}
(1) \unit{m^2} gedeeld door \unit{m^2/s} geeft seconden. (2) $400/10^{-6} =
\qty{4e8}{s}$, dat wil zeggen $4 \times 10^8 / 3.16 \times 10^7 \approx 12.7$
jaar. (3) $0.21 \times 12.7 \approx 2.7$ jaar.
\end{solution}

\begin{solution}{exo:g11:igneous-rocks:9}
Het is \omterm{def:g11:igneous-rocks:porphyritic}{porfirisch}: \omterm{def:g11:igneous-rocks:porphyritic}{fenokristen} van plagioklaas in een fijnkorrelige grondmassa.
Eerst koelde het magma langzaam af op diepte, en groeiden de plagioklaaskristallen
groot terwijl het grootste deel van het magma vloeibaar bleef; dan steeg het op en koelde het snel af,
en kristalliseerde de overgebleven vloeistof als de fijne grondmassa.
\end{solution}

\begin{solution}{exo:g11:igneous-rocks:10}
(1) Links: 2 Mg, 2 Si, 6 O; rechts: 2 Mg, 2 Si, 6 O. (2) \qty{1}{kg}
olivijn is $1000/140.7 \approx 7.11$~mol, dat evenveel mol
silica opneemt: $7.11 \times 60.1 \approx \qty{427}{g}$. (3) Als er geen silica meer in
de vloeistof is, kan de reactie niet doorgaan: het olivijn blijft.
\end{solution}

\begin{solution}{exo:g11:igneous-rocks:11}
(1) $2.5 \times 20 \times 1100 = \qty{55000}{MJ}$, \qty{5.5e10}{J}.
(2) $900 \times 20 = \qty{18000}{MJ}$, \qty{1.8e10}{J}. (3) De
kristallisatiewarmte is ongeveer een derde van de warmte die door afkoeling verloren gaat, een kwart van
het totaal: genoeg om het bevriezen sterk te vertragen, omdat ze helemaal
boven de solidus vrijkomt.
\end{solution}

\begin{solution}{exo:g11:igneous-rocks:12}
(1) De kleinste kristallen bij het contact, dat binnen een dag bevroor; de
grootste in het midden, dat ongeveer 950 dagen vloeibaar bleef. (2) Nee: de
eerste \qty{5}{m} duurden ongeveer 545 dagen, de laatste \qty{5}{m} ongeveer 400 dagen; het
front vertraagt eerst, en versnelt dan dicht bij het midden, waar de warmte die
aan beide kanten wegstroomt elkaar ontmoet. (3) Zelfs in zijn midden bevroor de \omterm{def:g11:igneous-rocks:intrusion}{sill} in
minder dan drie jaar; een graniet met grote kristallen heeft een veel langere tijd nodig.
\end{solution}

\begin{solution}{exo:g11:igneous-rocks:13}
(1) De \omterm{def:g11:igneous-rocks:intrusion}{pluton} is 100 keer dikker: $2.6 \times 100^2 = \num{26000}$
jaar. (2) Bijvoorbeeld: een diepe \omterm{def:g11:igneous-rocks:intrusion}{pluton} zit in gesteente dat al heet is, dus
verliest hij langzamer warmte dan aan gesteente van \qty{50}{\degreeCelsius}; hij kan
keer op keer door nieuw magma worden gevoed; een \omterm{def:g11:igneous-rocks:intrusion}{pluton} is geen oneindige plaat;
de eigenschappen van gesteente veranderen met de temperatuur.
\end{solution}

\begin{solution}{exo:g11:igneous-rocks:14}
Een \omterm{def:g11:igneous-rocks:felsic-mafic}{mafisch} magma is heet en arm aan silica: het kristalliseert de eerste mineralen
van de reeks --- olivijn, pyroxeen, calciumrijk plagioklaas --- en raakt door zijn
silica heen voordat er kwarts kan ontstaan. Een \omterm{def:g11:igneous-rocks:felsic-mafic}{felsisch} magma is koeler en rijk aan
silica: olivijn is er niet stabiel in, en het gaat door tot het einde van de
reeks, waarbij natriumrijk plagioklaas, kaliveldspaat en kwarts ontstaan.
\end{solution}

\begin{solution}{exo:g11:igneous-rocks:15}
(1) $b = \sqrt{t\kappa/C} = \sqrt{86400 \times 10^{-6}/0.21} \approx
\qty{0.64}{m}$. (2) De wet geeft $0.21 \times 25/10^{-6} \approx
\qty{5.3e6}{s}$, ongeveer twee maanden: dezelfde orde als een maand (een lavastroom
verliest ook warmte aan de lucht aan zijn bovenkant). (3) Glasachtig voor materiaal dat in
lucht of water wordt uitgeworpen; fijnkorrelig voor lavastromen en dunne platen; grover, maar nog
klein, voor dikke ondiepe \omterm{def:g11:igneous-rocks:intrusion}{intrusies}.
\end{solution}

\begin{solution}{pb:g11:igneous-rocks:1}
\textbf{1.} Een \omterm{def:g11:igneous-rocks:intrusion}{sill}: hij ligt evenwijdig aan de lagen, ertussen.

\textbf{2.} Bij het contact heel fijnkorrelig, bijna glasachtig; in het midden
fijn- tot middelkorrelig, de mineralen net zichtbaar met een loep.

\textbf{3.} Het contact verloor zijn warmte het snelst, aan de koude zandsteen: de
kristallen hadden geen tijd om te groeien.

\textbf{4.} Een afschrikrand. Hij bewijst dat het magma heet was toen het op zijn
plaats kwam, en de zandsteen koud.

\textbf{5.} Een \omterm{def:g11:igneous-rocks:intrusion}{gang}, de toevoer van de \omterm{def:g11:igneous-rocks:intrusion}{sill}.

\textbf{6.} \omterm{def:g11:igneous-rocks:felsic-mafic}{Mafisch}: \qty{49}{\percent} \ce{SiO2}, rijk aan ijzer, magnesium en
calcium.

\textbf{7.} Olivijn, pyroxeen en calciumrijk plagioklaas.

\textbf{8.} Een basalt; het plutonische gesteente van dezelfde samenstelling is \omterm{def:g11:ocean-ridges:gabbro}{gabbro}.

\textbf{9.} De \omterm{def:g11:igneous-rocks:texture}{textuur} hangt af van de afkoelingssnelheid, die verschilt van het
contact tot het midden, niet van de samenstelling.

\textbf{10.} Olivijn eerst; plagioklaas begint ongeveer op hetzelfde punt als
pyroxeen.

\textbf{11.} $b^2/\kappa = 400/10^{-6} = \qty{4e8}{s}$.

\textbf{12.} $0.206 \times 4 \times 10^8 \approx \qty{8.2e7}{s}$, dat wil zeggen
$8.2 \times 10^7 / 3.16 \times 10^7 \approx 2.6$ jaar.

\textbf{13.} $0.206 \times 4/10^{-6} \approx \qty{8.2e5}{s}$, ongeveer 9.5
dagen.

\textbf{14.} $0.07 \times 4 \times 10^8 = \qty{2.8e7}{s}$, ongeveer 0.9 jaar.
Kristalliserend magma geeft warmte vrij, die ook moet worden afgevoerd voordat
het magma vast wordt.

\textbf{15.} Ongeveer 950 dagen tegenover 0.6 dag: ongeveer 1600 keer sneller.

\textbf{16.} Ongeveer 950 dagen, zo'n 2.6 jaar: veel langer dan een lavastroom,
binnen de maanden tot jaren van ondiepe \omterm{def:g11:igneous-rocks:intrusion}{intrusies}.

\textbf{17.} Duizend keer langer vraagt een dikte die $\sqrt{1000} \approx
32$ keer groter is: ongeveer \qty{630}{m}.

\textbf{18.} Dicht bij het oppervlak zou een lichaam magma uitbarsten, of te snel afkoelen;
een groot lichaam dat heel lang afkoelt, moet diep liggen, in gesteente dat
al heet is.

\textbf{19.} Ze zijn door erosie weggesleten.

\textbf{20.} De \omterm{def:g11:igneous-rocks:intrusion}{sill} deed er ongeveer \textbf{2.6 jaar} over om te stollen.
\end{solution}
